\chapter{测量与周长}\label{ch:g4:measure}

三年级已经用尺子、秤和钟表做过测量
（见\cref{ch:g3:measure}）。本章增加单位换算表，它像一台能换算一切的“机器”，
以及测量平面图形的两种方法：图形边界的长度（\omterm{def:g4:measure:perimeter}{周长}），
和图形内部占据的大小（\omterm{def:g4:measure:area}{面积}）。
目前\omterm{def:g4:measure:area}{面积}先用方格数来计算。

\section{单位换算表}

\begin{method}[用换算表进行换算]\label{met:g4:measure:table}
对于长度，每一列都是其右边相邻一列的十倍：
\begin{center}
\renewcommand{\arraystretch}{1.3}
\begin{tabular}{c|c|c|c|c|c|c}
千米 & 百米 & 十米 & 米 & 分米 & 厘米 & 毫米 \\
\hline
 & & & $3$ & $4$ & $0$ & \\
\end{tabular}
\end{center}
\begin{enumerate}
  \item 每个数位占一列写出数字，测量数的个位数字
        写在它所用单位的那一列（例如 $3.4$ m 中的 $3$
        写在 m 的下面）；
  \item 按新单位读出这个数：$3.4$ m $= 340$ cm（空缺的
        数位用零补齐）；
  \item 质量（kg\,\dots\,g）和容量（L\,\dots\,cL）的换算方法
        与此相同。
\end{enumerate}
\end{method}

\begin{example}\label{ex:g4:measure:convert}
$5$ km $= 5\,000$ m；\quad $270$ cm $= 2.7$ m（$2$ 落在
m 的下面）；\quad $3$ kg $250$ g $= 3\,250$ g；\quad
$2.5$ L $= 250$ cL。
\end{example}

\section{周长}

\begin{definition}[周长]\label{def:g4:measure:perimeter}
\omterm{def:g4:geometry:polygon}{多边形}的\emph{周长}\index{周长}是其边界的总长度：
把所有边的长度相加（注意：各边长度的单位要相同！）。
常见图形的计算公式如下：
\[
\text{长方形：} P = 2 \times (L + w),
\qquad
\text{正方形：} P = 4 \times c .
\]
\end{definition}

\begin{example}\label{ex:g4:measure:perimeter}
一个长方形花园，长 $12$ m，宽 $7$ m：
\[
P = 2 \times (12 + 7) = 2 \times 19 = 38 \text{ m}
\]
所用的篱笆长度。一个边长为 $18$ cm 的正方形相框：
$P = 4 \times 18 = 72$ cm。
\end{example}

\section{面积：数方格}

\begin{definition}[数方格求面积]\label{def:g4:measure:area}
画在方格纸上的图形的\emph{面积}\index{面积}，是它所覆盖的
方格的个数。两个半格合起来算作一格。面积
回答的是“图形有多大”，\omterm{def:g4:measure:perimeter}{周长}回答的是“边界有多长”——
这是两个不同的问题！
\end{definition}

\begin{omfigure}
\begin{tikzpicture}[scale=0.55]
  \draw[gray!40, thin] (0,0) grid (13,4);
  \draw[very thick, omDef, fill=omDef!15] (1,1) rectangle (5,3);
  \node[font=\small, omDef] at (3,2) {$8$ 个方格};
  \draw[very thick, omThm, fill=omThm!15]
    (7,1) -- (11,1) -- (11,2) -- (9,2) -- (9,3) -- (7,3) -- cycle;
  \node[font=\small, omThm] at (8.4,1.5) {$6$ 个方格};
\end{tikzpicture}

{\small 方格纸上的两个图形：\omterm{def:g4:measure:area}{面积}分别为 $8$ 和 $6$ 个方格。现在沿边界走一圈，数一数
单位长度的个数：两者都是 $12$！\omterm{def:g4:measure:perimeter}{周长}相同，
\omterm{def:g4:measure:area}{面积}却不同。}
\end{omfigure}

\begin{example}[半格]\label{ex:g4:measure:half}
方格纸上画出一个\omterm{def:g3:shapes:rightangle}{直角}三角形，两条直角边分别为 $4$ 和 $2$，它覆盖了 $3$ 个整格
和 $2$ 个半格\,\dots\ 数的时候请仔细：它的\omterm{def:g4:measure:area}{面积}
是 $4 \times 2$ 长方形\omterm{def:g4:measure:area}{面积}的\omterm{def:g3:division:half}{一半}，即 $4$ 个方格。沿长方形的对角线剪开，
总会得到
\omterm{def:g4:measure:area}{面积}相等的两个三角形。
\end{example}

\begin{example}[面积相同，周长不同]\label{ex:g4:measure:samearea}
一个 $4 \times 4$ 的正方形和一个 $8 \times 2$ 的长方形都覆盖 $16$
个方格，\omterm{def:g4:measure:area}{面积}相同。边界：正方形为 $16$ 个单位，长方形为 $20$
个单位。两者互不决定：狭长的图形边界很长，
而\omterm{def:g4:measure:area}{面积}却很小。
\end{example}

\section{时长}

\begin{method}[跨过整时的时间加减]\label{met:g4:measure:time}
做时长的加减时，以整时为节点分步进行：
（回忆：$1$ h $= 60$ min）：
\[
9\text{ h }40 + 35 \text{ min}:
\quad 9\text{ h }40 \xrightarrow{+20} 10\text{ h }00
\xrightarrow{+15} 10\text{ h }15 .
\]
切勿不先换算就把分钟数加过 $60$：$40 + 35 = 75$ min
$= 1$ h $15$ min。
\end{method}

\begin{example}\label{ex:g4:measure:timetable}
火车在 $8{:}47$ 出发，$11{:}05$ 到达。时长：
$8{:}47 \to 9{:}00$ 为 $13$ min；$9{:}00 \to 11{:}00$ 为 $2$ h；
$11{:}00 \to 11{:}05$ 为 $5$ min。合计：$2$ h $18$ min。
\end{example}

\section{练习}

\begin{exercise}[$\star$]\label{exo:g4:measure:1}
利用单位换算表换算：$7$ m 化为 cm；$3.2$ km 化为 m；$85$ mm 化为
cm；$4\,500$ m 化为 km。
\end{exercise}

\begin{exercise}[$\star$]\label{exo:g4:measure:2}
换算：$2$ kg 化为 g；$1\,800$ g 化为 kg 和 g；$3.5$ L 化为 cL；
$25$ cL 化为 L。
\end{exercise}

\begin{exercise}[$\star$]\label{exo:g4:measure:3}
计算下列图形的\omterm{def:g4:measure:perimeter}{周长}：长方形，$9$ cm $\times$ $5$ cm；边长为
$7.5$ cm 的正方形；三角形，三边长分别为 $6$ cm、$8$ cm 和 $10$ cm。
\end{exercise}

\begin{exercise}[$\star$]\label{exo:g4:measure:4}
一个长方形的\omterm{def:g4:measure:perimeter}{周长}为 $26$ cm，长为 $8$ cm。求它的宽
（先求长 $+$ 宽）。
\end{exercise}

\begin{exercise}[$\star$]\label{exo:g4:measure:5}
在方格纸上画一个 L 形图形，恰好覆盖 $10$
个方格，并数出它的\omterm{def:g4:measure:perimeter}{周长}包含几个单位长度。
\end{exercise}

\begin{exercise}[$\star$]\label{exo:g4:measure:6}
画出两个\omterm{def:g4:measure:area}{面积}都为 $12$ 个方格的不同长方形。计算每个长方形的
\omterm{def:g4:measure:perimeter}{周长}。哪一个更“圆”，哪一个更“细”？
\end{exercise}

\begin{exercise}[$\star$]\label{exo:g4:measure:7}
在方格纸上画一个\omterm{def:g3:shapes:rightangle}{直角}三角形，两直角边长分别为 $6$ 和 $3$。它的\omterm{def:g4:measure:area}{面积}
是多少个方格？（用\cref{ex:g4:measure:half}。）
\end{exercise}

\begin{exercise}[$\star$]\label{exo:g4:measure:8}
计算下列时长：从 $10{:}20$ 到 $12{:}45$；从 $7{:}35$ 到
$9{:}10$。一部电影时长 $1$ h $50$，从 $20{:}30$ 开始播放，
问它何时结束？
\end{exercise}

\begin{exercise}[$\star$]\label{exo:g4:measure:9}
基姆绕一块长方形场地跑 $3$ 圈，场地长 $60$ m、宽 $45$ m。
她共跑了多少米？
\end{exercise}

\begin{exercise}[$\star\star$]\label{exo:g4:measure:10}
一位农夫想给边长为 $85$ m 的正方形田地围篱笆，
留出一个 $4$ m 的门口不围篱笆。他需要买多少米篱笆？
\end{exercise}

\begin{exercise}[$\star\star$]\label{exo:g4:measure:11}
判断下列说法的对错，并用方格图举例：“若两个图形的
\omterm{def:g4:measure:perimeter}{周长}相同，则它们的\omterm{def:g4:measure:area}{面积}相同”；“若两个图形的
\omterm{def:g4:measure:area}{面积}相同，则它们的\omterm{def:g4:measure:perimeter}{周长}相同”。
（可参考 \Cref{ex:g4:measure:samearea} 及本章的图。）
\end{exercise}
